% Pregunta 617 (2009-octubre-OEP2009, n.º 14): subida de r* con movilidad perfecta de capitales y tipo de cambio flexible; a la izquierda, precios rígidos (Mundell-Fleming en el plano Y-epsilon); a la derecha, precios flexibles (NX(epsilon) = S - I(r*)). Misma escala vertical: la depreciación real es mayor con precios rígidos. Autor: Víctor Gutiérrez Marcos
\documentclass[tikz,border=6pt]{standalone}
\input{tikz-tcee}
\usetikzlibrary{decorations.pathreplacing,calc}
\begin{document}
\begin{tikzpicture}[tcee,x=1cm,y=1cm]
  % ---------- Panel 1: precios rígidos ----------
  % IS*: e = 4,6 - 0,6Y ; IS*': e = 3,40 - 0,6Y (baja |DI| = 1,2, como en el panel 2) ; LM*: Y0 = 2,2 ; LM*': Y1 = 3,6
  % E0: (2,2 ; 3,28) ; E1: (3,6 ; 1,24)  -> caída de e = 2,04 = 0,6·1,4 + 1,2 (con NX de pendiente -1 como en el panel 2)
  \node[font=\small\bfseries] at (3,5.5) {Precios rígidos};
  \ejes{$Y$}{$\varepsilon$}{6}{5}
  \draw[curva1] (2.2,0.2) -- (2.2,4.7) node[etiqueta,above] {$LM^*$};
  \draw[curvanueva,desplazada] (3.6,0.2) -- (3.6,4.7) node[etiqueta,above] {$LM^{*\prime}$};
  \draw[curva2] (0.4,4.36) -- (5.4,1.36) node[etiqueta,right] {$IS^*$};
  \draw[curvanueva,desplazada] (0.4,3.16) -- (5.2,0.28) node[etiqueta,right] {$IS^{*\prime}$};
  \draw[guia] (2.2,3.28) -- (0,3.28) node[left] {$\varepsilon_0$};
  \draw[guia] (3.6,1.24) -- (0,1.24) node[left] {$\varepsilon_1$};
  \node[below] at (2.2,0) {$Y_0$};
  \node[below] at (3.6,0) {$Y_1$};
  \node[punto] at (2.2,3.28) {};
  \node[punto,tceerojo] at (3.6,1.24) {};
  \node[above right] at (2.2,3.28) {$E_0$};
  \node[above right] at (3.6,1.24) {$E_1$};
  \draw[flecha] (2.75,4.75) -- (3.35,4.75);
  \draw[flecha] (1.15,3.95) -- (0.75,3.4);
  \draw[decorate,decoration={brace,amplitude=4pt,mirror},thick,tceerojo] (-0.55,3.28) -- (-0.55,1.24) node[midway,left=5pt,text=tceerojo] {$\Delta\varepsilon$};

  % ---------- Panel 2: precios flexibles ----------
  % NX: e = 4,78 - NX ; S - I0 en NX = 1,5 ; S - I1 en NX = 2,7 (|DI| = 1,2)
  % E0: (1,5 ; 3,28) ; E1: (2,7 ; 2,08)  -> caída de e = 1,20
  \begin{scope}[xshift=8.6cm]
  \node[font=\small\bfseries] at (3,5.5) {Precios flexibles ($Y=\overline{Y}$)};
  \ejes{$NX$}{$\varepsilon$}{6}{5}
  \draw[curva1] (0.3,4.48) -- (4.5,0.28) node[etiqueta,right] {$NX(\varepsilon)$};
  \draw[curva2] (1.5,0.2) -- (1.5,4.7) node[etiqueta,above] {$S-I(r^*_0)$};
  \draw[curvanueva,desplazada] (2.7,0.2) -- (2.7,4.7) node[etiqueta,above right=0pt and -6pt] {$S-I(r^*_1)$};
  \draw[guia] (1.5,3.28) -- (0,3.28) node[left] {$\varepsilon_0$};
  \draw[guia] (2.7,2.08) -- (0,2.08) node[left] {$\varepsilon_1$};
  \node[punto] at (1.5,3.28) {};
  \node[punto,tceerojo] at (2.7,2.08) {};
  \node[above right] at (1.5,3.28) {$E_0$};
  \node[above right] at (2.7,2.08) {$E_1$};
  \draw[decorate,decoration={brace,amplitude=4pt},thick] (1.5,-0.55) -- (2.7,-0.55) ;
  \node[below] at (2.1,-0.62) {$|\Delta I|$};
  \draw[flecha] (1.7,4.2) -- (2.5,4.2);
  \draw[decorate,decoration={brace,amplitude=4pt,mirror},thick,tceerojo] (-0.55,3.28) -- (-0.55,2.08) node[midway,left=5pt,text=tceerojo] {$\Delta\varepsilon$};
  \end{scope}
\end{tikzpicture}
\end{document}
